% ECONOMICS_PROBLEM: {"id": "C73-15", "title": "Payoff characterization of convergence for a fixed mirror-descent rule", "jel_code": "C73", "source_id": "OP-0148"}
% !TeX program = xelatex
\documentclass[11pt,a4paper]{article}
\usepackage[margin=23mm]{geometry}
\usepackage{fontspec}
\setmainfont{Latin Modern Roman}
\usepackage{amsmath,amssymb,mathtools}
\usepackage{xcolor,enumitem,booktabs,longtable,array,xurl,hyperref}
\definecolor{muted}{HTML}{53636C}
\hypersetup{colorlinks=true,urlcolor=blue,linkcolor=blue}
\setlength{\parindent}{0pt}
\setlength{\parskip}{5pt}
\newcommand{\R}{\mathbb R}
\newcommand{\E}{\mathbb E}
\newcommand{\N}{\mathbb N}
\newcommand{\ind}{\mathbf 1}
\newcommand{\Var}{\operatorname{Var}}
\newcommand{\Cov}{\operatorname{Cov}}
\newcommand{\supp}{\operatorname{supp}}
\DeclareMathOperator*{\argmin}{arg\,min}
\DeclareMathOperator*{\argmax}{arg\,max}
\newcommand{\entrymeta}[3]{{\sffamily\small\color{muted}\textbf{\texttt{#1}}\quad|\quad #2\par #3\par}}
\newcommand{\Problemref}[1]{\texttt{#1}}
\newcommand{\JELref}[2]{\texttt{#2} (JEL \texttt{#1})}
\begin{document}
\section*{Payoff characterization of convergence for a fixed mirror-descent rule}
\textbf{Problem ID:} \texttt{C73-15}\quad \textbf{JEL:} \texttt{C73}\par
% BEGIN ECONOMICS STATEMENT
\label{problem:OP-0148}
% GAME_THEORY_PROBLEM: {"local_id": "GT-018", "original_number": 18, "kind": "R", "entry_kind": "statement_only", "published": false, "review_required": true, "category": "game_theory"}
\label{game:GT-018}
{\sffamily\small\color{muted}
\textbf{GT-018} | Learning dynamics and potential structure\\
\textbf{R} | Formal algorithm-specific classification agenda.
\par}
\subsection*{Definitions and mathematical statement}
Fix action sets $S_i$ and proper lower-semicontinuous strictly convex regularizers $h_i:\Delta(S_i)\to\mathbb R\cup\{+\infty\}$, finite and differentiable on the relative interiors, together with a positive step-size sequence $(\eta_t)$. Let
\[
 D_{h_i}(p,x)=h_i(p)-h_i(x)-\langle\nabla h_i(x),p-x\rangle
\]
be the Bregman divergence. In a finite game $G$, let $v_i(x)$ be the vector with coordinates $u_i(a_i,x_{-i})$. A specified simultaneous online mirror-descent rule is
\[
 x_i^{t+1}=\underset{p\in\Delta(S_i)}{\operatorname{argmax}}
 \{\eta_t\langle v_i(x^t),p\rangle-D_{h_i}(p,x_i^t)\}.
\]
Assume the declared regularizers and step sizes make each update uniquely defined and remain within the domain of $\nabla h_i$. For this fixed rule $D$, define
\[
 \mathcal G_D=\{G:\ \forall x^1\in\operatorname{relint}X,
 \ \lim_{t\to\infty}\operatorname{dist}(x^t,\operatorname{NE}(G))=0\}.
\]
Give necessary and sufficient conditions, in terms of the game's payoff arrays, for $G\in\mathcal G_D$.
\subsection*{Short English statement}
Which games make a specified adaptive update converge to the Nash set from every interior starting point?
\subsection*{Variants and scope boundaries}
The update rule, learning rates, and convergence notion cannot be left implicit. Potential structure does not guarantee convergence for every arbitrary discrete-time step size. Statistical prediction of convergence from a scalar potentialness measure is different from a necessary-and-sufficient theorem. Requiring convergence to a single equilibrium strengthens set convergence.
\subsection*{Source relationship and status}
[S1] studies potentialness and explicitly calls the general characterization problem open. The displayed mirror-descent family, universal initialization requirement, and fixed schedule are editorial choices that make one member of that agenda a defined mathematical target.
\subsection*{References}
\begingroup\footnotesize\raggedright
{\footnotesize\raggedright
\textbf{[S1]} Martin Bichler, Davide Legacci, Panayotis Mertikopoulos, Matthias Oberlechner, and Bary Pradelski (2025). \emph{Characterizing the Convergence of Game Dynamics via Potentialness}. arXiv:2503.16285v1. Locator: Sections 3--4 for games and potentialness; Section 6 for the convergence-classification agenda. \url{https://arxiv.org/html/2503.16285v1}\par}
\par\endgroup

\subsection*{Common conventions: Source mathematical conventions}

\textbf{Finite games.} For a finite set $A$, $\Delta(A)$ is its probability
simplex. Players choose independent mixed strategies in Nash equilibrium;
correlation is allowed only when explicitly part of the solution concept.
In a game with payoff functions $u_i$ and strategy profile $x$, an additive
$\varepsilon$ Nash equilibrium satisfies
\[
 u_i(x)\geq u_i(x_i',x_{-i})-\varepsilon
 \quad\text{for every player $i$ and every admissible deviation $x_i'$}.
\]
For numerical approximation, payoffs are normalized as locally specified.
An $\varepsilon$ payoff residual is distinct from distance at most
$\varepsilon$ to the set of exact equilibria. Point convergence, convergence
of empirical play, and convergence in distance to an equilibrium set are
also distinct.

\textbf{Computational representations.} Unless stated otherwise, a complexity
question takes rational data with binary encoding; polynomial time is measured
in total bit length. A fixed-accuracy polynomial algorithm is not necessarily
polynomial in $1/\varepsilon$, and neither is necessarily polynomial in
$\log(1/\varepsilon)$. Succinct games and value, demand, or sampling oracles
must declare their interfaces. Exact real solutions require an explicit
representation; an unspecified list of arbitrary real numbers is not a
finite computational output. A claimed algorithmic barrier is meaningful
only for its specified model and reductions.

\textbf{Dynamic games.} Histories, information, observation, and allowable
randomization are part of the model. A stationary strategy depends only on
the current state; a positional strategy in a deterministic graph game
chooses one action at each controlled vertex. Nash equilibrium at an initial
state and equilibrium after every history are different requirements.
An expectation of a pathwise $\liminf$ payoff, a $\liminf$ of expected averages,
a limiting expected average, and a uniform finite-horizon equilibrium are
different criteria. None is silently substituted for another.

\textbf{Allocation and mechanisms.} A complete allocation assigns every item
exactly once unless otherwise stated. Zero-valued goods, negative values,
capacity constraints, feasibility, lotteries, and copies of goods affect
fairness definitions and are specified locally. Dominant-strategy incentive
compatibility, Bayesian incentive compatibility, universal truthfulness,
and truthfulness in expectation impose different inequalities. Whenever
an objective ratio has zero denominator, its convention is stated or those
instances are excluded.

\textbf{Infinite-dimensional models.} Expectations, integrals, stochastic
equations, and optimization objectives require their local regularity,
integrability, filtrations, and admissible domains. A backward--forward
mean-field system is a boundary-value problem; it is not an initial-value
semiflow without an additional construction. Long-time, large-population,
and vanishing-noise limits require an order of limits and a convergence
topology. If a minimum need not be attained, the objective is its infimum
or supremum and the approximation requirement is stated separately.
% END ECONOMICS STATEMENT
\end{document}
